\documentclass[aspectratio=169,10pt]{beamer}
\usetheme{metropolis}
\usepackage{booktabs,amsmath,amssymb}
\title{Q10: The Jagged Frontier}
\subtitle{Why 'simple' is a human notion the model never learned}
\author{DLS 2026 --- Kukla}
\date{October 2026}
\begin{document}
\maketitle

\begin{frame}{The puzzle}
\begin{block}{Student question (\#18)}
Why can models solve complex problems yet still make simple mistakes?
\end{block}
It explains relativity, then miscounts letters in ``strawberry''.
\end{frame}
\begin{frame}{Human vs model competence}
\begin{itemize}
\item Humans: competence is \emph{correlated} (hard $\Rightarrow$ easy follows).
\item Models: a \textbf{jagged frontier} --- sharp, unpredictable variation.
\end{itemize}
\alert{``Hard'' and ``easy'' are human notions the model never learned.}
\end{frame}
\begin{frame}{Root cause 1: the objective}
\begin{itemize}
\item Trained to predict next token, not to be uniformly competent.
\item Capability tracks \textbf{training frequency}, not logical difficulty.
\item Relativity prose is abundant; exact long arithmetic is rare.
\end{itemize}
\end{frame}
\begin{frame}{Root cause 2: tokenization}
\begin{itemize}
\item Model sees sub-word tokens, not letters or digits.
\item Counting r's needs letter structure it doesn't represent.
\item Number chunks break positional arithmetic.
\end{itemize}
\alert{Predictable representation artifacts, not ``stupidity''.}
\end{frame}
\begin{frame}{Root cause 3: no built-in checking}
\begin{itemize}
\item One forward pass, no scratchpad, no verification.
\item Humans re-check; a base model doesn't unless prompted / given a tool.
\item Same model often fixes it the instant you say ``think step by step''.
\end{itemize}
\end{frame}
\begin{frame}{Predict the jaggedness}
Strong where: well-represented, error-tolerant, pattern-like.\\
Weak where: exact symbols, long chains, sub-token info, rare knowledge.\\[1em]
\alert{Pair brittle regimes with reasoning, tools, and verifiers.}
\end{frame}

\end{document}
